\section{Заключение}

Работа оценивает предельную информационную ценность распределительных статистик домохозяйств для
правила процентной ставки при неполной наблюдаемости в локальной HANK/SSJ-среде. В текущей версии
основной технической спецификацией является совместный фильтр Калмана: агрегатные и
распределительные оценки строятся из одной линейной модели состояния, а информационные режимы
различаются только матрицей наблюдения.

Отдельная конечная разностная проверка поддерживает локальную SSJ-аппроксимацию в используемом
диапазоне шоков. Экспортированный якобиан денежного шока даёт малую среднюю ошибку относительно
нелинейного HANK-перехода, а неполитические шоки также демонстрируют небольшие ошибки локальной
линейности. Это не превращает расчёт в глобальную нелинейную HANK-оценку, но делает локальную
интерпретацию результатов существенно более защищённой.

Главный устойчивый результат состоит в том, что восстановление агрегатного состояния существенно
снижает потери относительно текущих шумных агрегатов. Распределительная информация даёт
дополнительный, но умеренный выигрыш сверх фильтрованных агрегатов в базовой парной оценке и в
ряде идентификационных проверок. Однако large-sample проверка с отдельными HANK/SSJ-шоками и
кластеризацией по макрошоку делает этот вывод осторожнее: полный распределительный блок уже не
имеет статистически устойчивого выигрыша сверх фильтрованных агрегатов, тогда как отдельные
распределительные статистики сохраняют положительный, но слабый знак.

Идентификационные проверки показывают, что этот результат не сводится к механическому увеличению
числа признаков. Искусственные признаки с похожей дисперсией, авторегрессией и корреляцией с
агрегатами не воспроизводят выигрыш; временно сдвинутые признаки ухудшают правило. При этом
остаточная распределительная информация, очищенная от линейной зависимости от фильтрованных
агрегатов, сохраняет положительный эффект в базовом тесте. Более аккуратный вывод таков:
распределительные статистики могут быть самостоятельным источником информации для правила ставки,
если они помогают оценивать будущую трансмиссию ставки сверх уже восстановленного агрегатного
состояния, но этот канал в текущей калибровке не является настолько сильным, чтобы полный
распределительный блок устойчиво проходил large-sample кластерную проверку. Это именно локальный
результат HANK/SSJ-эксперимента, а не решение глобальной нелинейной задачи оптимальной политики в
полной HANK-модели.

Проверка с обратной связью усиливает этот вывод для локальной SSJ-постановки. После добавления
прямых HANK-откликов распределительных статистик на траекторию ставки полное распределительное
состояние устойчиво снижает потери относительно фильтрованных агрегатов. При этом результат
остаётся локальным: это не глобальная нелинейная HANK-оптимизация, а проверка замороженных правил в
локальной проекции вокруг стационарного состояния.

Поэтому финальная позиция работы не состоит в том, что распределительная информация всегда нужна
для денежно-кредитной политики. Более защищаемый вывод другой: в локальной HANK/SSJ-среде
агрегатная фильтрация является устойчиво ценной, а распределительные статистики имеют потенциальную
предельную ценность в тех спецификациях, где они несут информацию о будущей трансмиссии ставки.
Наиболее сильное подтверждение этого канала получается в closed-loop проверке с прямыми
распределительными SSJ-откликами; наиболее осторожный результат -- в large-sample проверке
открытой локальной проекции, где полный распределительный блок не проходит кластерную проверку.
